\specialsection{函数空间}

设 G 是一个局部紧交换群。我们回顾 G 上的主要函数空间以及它们与 Fourier 变换的关系。

左列定义 G 上的一个函数空间 E，并指出它是 \(L^{1}(G)\) 的子空间、\(L^{2}(G)\) 的子空间，还是两者的子空间。右列对 E 在 Fourier 变换和余变换下的象指出类似的包含关系（后者依情形或是定义在 \(L^{1}(G)\) 上的那个，或是定义在 \(L^{2}(G)\) 上的那个）。记号 F 表示 Fourier 变换或 Fourier 余变换。

\(\mathcal{M}^1 (\mathrm{G})\)

\[
\mathcal {F} (\mathcal {M} ^ {1} (\mathrm{G})) \subset \mathcal {C} _ {b} (\widehat {\mathrm{G}})
\]

（有界测度的空间）

（命题 3，第 207 页）

\(\mathrm{L}^{1}(\mathrm{G})\)

\[
\mathcal {F} (\mathrm{L} ^ {1} (\mathrm{G})) \subset \mathcal {C} _ {0} (\widehat {\mathrm{G}})
\]

（可积函数）

\(\mathrm{L}^2 (\mathrm{G})\)

\[
\mathscr {F} (\mathrm{L} ^ {2} (\mathrm{G})) = \mathrm{L} ^ {2} (\widehat {\mathrm{G}})
\]

（平方可积函数）（定理 1，第 215 页）

\(\mathrm{B(G)}\subset \mathrm{L}^{1}(\mathrm{G})\) \(\mathcal{F}(\mathrm{B(G)}) = \mathrm{B}(\widehat{\mathrm{G}})\)

由满足如下条件的 \(f \in \mathrm{L}^1(\mathrm{G})\) 组成的空间（定理 3，第 222 页）

\[
\mathcal {F} _ {\mathrm{G}} (f) \in \mathrm{L} ^ {1} (\widehat {\mathrm{G}})
\]

\(\mathrm{A(G)}\subset \mathrm{L}^1 (\mathrm{G})\cap \mathrm{L}^2 (\mathrm{G})\cap \mathrm{B(G)}\) \(\mathcal{F}(\mathrm{A(G)})\subset \mathcal{C}_0(\widehat{\mathrm{G}})\cap \mathrm{L}^1 (\widehat{\mathrm{G}})\cap \mathrm{L}^2 (\widehat{\mathrm{G}})\) （由 \(f*g\) 张成的空间，其中 \(f\)、\(g\in \mathrm{L}^{1}(\mathrm{G})\cap \mathrm{L}^{2}(\mathrm{G}))\)）（命题 8 的推论，第 211 页，引理 5，第 215 页，命题 11，第 217 页）

